\paragraph{\textcolor{red}{Acknowledgment (remember to edit this).}} Our team members are \textbf{Isaac Newton (in123)}, \textbf{Albert Einstein (ae123)}, and \textbf{Alan Turing (at123)}. We searched for Problem 17 on the Internet after spending 123 hours on it, and we asked ChatGPT how to solve Problem 32(c) after spending 123 days on it. We did not successfully solve Problem 49 despite having spent 123 years on it. \textbf{We affirm that we have acknowledged every external help during the preparation of this document.}


\section*{Problem 1}
There are $n$ children forming a circle. The $i$-th child currently has $a_i$ apples. Let $\bar a = \frac{\sum_{i = 1}^n a_i}{n}$ be the average number of apples that a child has, and we assume that it is an integer.

In one move, we can take one apple from a child and give it to an \emph{adjacent} child. What is the minimum number of moves we have to make to make every child have exactly $\bar a$ apples? Model this problem by a flow network with costs, and therefore show that this problem can be solved in polynomial time.

\paragraph{\textcolor{brown}{Solution.}}
Using flow networks, let's create a directed graph where each child i is a vertex. We connect each child to neighboring children in a circle being i - 1 mod and i + 1 mod.

Three cases for supply/demand: 
    If a_i < a then i has demand a - a_i
    If a_i > a then i has supply a_i - a
    If neutral then a = a_i

This should give a $O(n)$ run-time as it is based on similar cost problems where the capacities and costs are bounded by polynomials.



\section*{Problem 2}
Explicitly construct a bijection between each of the following pairs of sets and therefore show that they have the same cardinality.
\begin{enumerate}
    \item Natural numbers and perfect squares.
    \item $(0, 1)$ and $(1, +\infty)$.
    \item $(-1, 1)$ and $\mathbb{R}$.
    \item $[0, 1]$ and $[0, 1)$.
\end{enumerate}


\paragraph{\textcolor{brown}{Solution.}}
\begin{enumerate}
    \item Define f : N -> {n^2 | n in the set N} as f(n) = n^2
            The function is injective because f(n1) = f(n2) would mean that n1^2 = n2^2 and that means n1 = n2 which is in the set of N
            The function is surjective because for any perfect sqauare m^2, there exists a natural number m such that f(m) = m^2
            This is a bijection
    \item   Define f : (0,1) -> (1, +infinity) as f(x) = 1/x
            The function is injective as if f(x1) = f(x2) then 1/x1 = 1/x2 implies x1 = x2
            The function is surjective as y is in the set of (1, +infinity) there exists x = 1/y in the set of (0,1) such that f(x) = y 
            This is a bijection
    \item   Define f : (-1,1) -> R as f(x) = x/1-|x|
            The function is injective as if f(x1) = f(x2) then x1/1-|x1| = x2/1 - |x2| which implies x1 = x2
            The function is surjective as for any y in the set of R there exists x = y/1 + |y| in the set (-1,1) such that f(x) = y
            This is a bijection
    \item   Define f : [0,1) as:
                f(x) = x if x doesn't equal 1/2^n for any n in the set of N
                f(1/2^n) = 1/2^n + 1 for n in the set of N
                f(0) = 0
                The function is injective as each element in [0,1) is mapped to by most one element
                The function is surjective because every element in [0,1) is mapped by some element
                This is a bijection
\end{enumerate}



\section*{Problem 3}
Review the halting problem, and prove that the following problem is also undecidable: Given a Turing machine, decide whether it always correctly decides the set cover problem. (Note that the given Turing machine may not halt on a valid input, in which case it does not correctly decide the set cover problem.)

\paragraph{\textcolor{brown}{Solution.}}
Here is my solution.


\section*{Problem 4}
Prove that the following problem is $\mathtt{NP}$-complete: Given $n$ positive integers that sum to $2W$, decide whether we can partition the integers into two parts with equal sum ($W$ each).

\paragraph{\textcolor{brown}{Solution.}}
Here is my solution.


\section*{Problem 5}
Prove that the following problem is $\mathtt{NP}$-complete: Given a weighted directed graph with $n$ vertices and $m$ edges, decide whether the graph has a simple cycle with sum of weights equal to $344$. The weight of each edge must be an integer, but can be positive, negative, or zero.

\paragraph{\textcolor{brown}{Solution.}}
Here is my solution.
